\documentclass{article}
\usepackage{graphicx} % Required for inserting images
\usepackage{amsmath}

\title{Seth Conjecture}
\author{Seth Berry}
\date{16 August 2026}

\begin{document}

\maketitle

\section{Introduction}

The Seth Conjecture needs the Seth function. So, what is that function, 
and what is this conjecture? It will be explained in the following sections.

\section{Pascal's triangle}

\begin{center}
\begin{tabular}{c}
1 \\
1 \quad 1 \\
1 \quad 2 \quad 1 \\
1 \quad 3 \quad 3 \quad 1 \\
1 \quad 4 \quad 6 \quad 4 \quad 1 \\
1 \quad 5 \quad 10 \quad 10 \quad 5 \quad 1 \\
1 \quad 6 \quad 15 \quad 20 \quad 15 \quad 6 \quad 1 \\
1 \quad 7 \quad 21 \quad 35 \quad 35 \quad 21 \quad 7 \quad 1 \\
1 \quad 8 \quad 28 \quad 56 \quad 70 \quad 56 \quad 28 \quad 8 \quad 1 \\
1 \quad 9 \quad 36 \quad 84 \quad 126 \quad 126 \quad 84 \quad 36 \quad 9 \quad 1 \\
1 \quad 10 \quad 45 \quad 120 \quad 210 \quad 252 \quad 210 \quad 120 \quad 45 \quad 10 \quad 1
\end{tabular}
\end{center}

To understand the Seth function, we first need to look at the rows of 
Pascal's triangle with the row function. The way the function works is 
it takes the diagonal row from left to right, e.g.,

$R_1(x) = 1, 1, 1, 1, 1, 1, \dots$ as the pattern, so $R_1(x) = x^0$

$R_2(x) = 1, 2, 3, 4, 5, 6, \dots$ as the pattern, so $R_2(x) = x$

$R_3(x) = 1, 3, 6, 10, 15, \dots$ as the pattern, so $R_3(x) = \frac{x(x + 1)}{2}$

and so on. This is the rows of the triangle. Now what is $R_x(x)$?

\section{Solving for $R_x(x)$ in terms of $n$}

To do this, let's make a formula for $R_n(x)$ so it is in terms of 
$x$ and $n$, then plug in $x$ for $n$ later. Now we need to find the pattern by starting with $R_4(x)$.

To start solving $R_4(x)$, we need to list the numbers in the function:

$R_4(x) = 1, 4, 10, 20, 35, 56, \dots$

Now we must look at the difference:
$3, 6, 10, 15, 21, \dots$

These are the triangular numbers from the previous row. Since the next number is found by adding the triangular number, we need to solve for:
$$\sum_{n=1}^{x} \frac{n(n + 1)}{2}$$

Now what we need to do is bring out the $\frac{1}{2}$:
$$\frac{1}{2} \sum_{n=1}^{x} n(n + 1)$$

Then we distribute inside the brackets to get $(n^2 + n)$ still inside the sigma:
$$\frac{1}{2} \sum_{n=1}^{x} (n^2 + n)$$

Then we pull them out into separate sigmas:
$$\frac{1}{2} \left( \sum_{n=1}^{x} n^2 + \sum_{n=1}^{x} n \right)$$

Now apply the formulas, then show the result of the formulas:
$$\frac{1}{2} \left[ \frac{x(x+1)(2x+1)}{6} + \frac{x(x+1)}{2} \right]$$

Now make the denominators the same:
$$\frac{1}{2} \left[ \frac{x(x+1)(2x+1)}{6} + \frac{3x(x+1)}{6} \right]$$

Then make it one fraction:
$$\frac{1}{2} \left[ \frac{x(x+1)(2x+1) + 3x(x+1)}{6} \right]$$

Then factor out $x(x+1)$, then do that:
$$\frac{1}{2} \left[ \frac{x(x+1)(2x+4)}{6} \right]$$

And then multiply the $(2x+4)$ by $\frac{1}{2}$ to get the final answer:
$$R_4(x) = \framebox[4cm]{$\displaystyle\frac{x(x+1)(x+2)}{6}$}$$

Now let's solve for $R_5(x)$. If we look at the previous differences, we 
see they are $1, 4, 10, 20, 35, \dots$, and looking at this, we see it 
is the same as the tetrahedral numbers or $R_4(x)$. So we must solve for:
$$\sum_{n=1}^{x} \frac{n(n+1)(n+2)}{6}$$

Now, using similar methods, we get:
$$R_5(x) = \framebox[4.8cm]{$\displaystyle\frac{x(x+1)(x+2)(x+3)}{24}$}$$

Then for $R_6(x)$:
$$R_6(x) = \framebox[5.2cm]{$\displaystyle\frac{x(x+1)(x+2)(x+3)(x+4)}{120}$}$$

Based on this, we can see that the answer for:
$$R_n(x) = \frac{x(x+1)(x+2)(x+3)\dots(x+n-2)}{(n-1)!}$$

Or, we can rewrite it by replacing $(n-1)!$ with $\Gamma(n)$:
$$R_n(x) = \framebox[6.5cm]{$\displaystyle\frac{x(x+1)(x+2)(x+3)\dots(x+n-2)}{\Gamma(n)}$}$$

At the beginning, we wanted to solve for $R_x(x)$, so our final answer is:
$$R_x(x) = \framebox[6.5cm]{$\displaystyle\frac{x(x+1)(x+2)(x+3)\dots(2x-2)}{\Gamma(x)}$}$$

\section{The Different Ways to Write $R_x(x)$}

There are multiple ways to rewrite $R_x(x)$. The first way is to see that we are always looking for the differences between the numbers, which means it is always based on the previous row. Therefore, we can write it in nested sigma form:
$$\sum_{a_1=1}^{x} \sum_{a_2=1}^{a_1} \sum_{a_3=1}^{a_2} \dots \sum_{a_x=1}^{a_{x-1}} \left( \lim_{b \to \infty} a_x^{-b} \right)$$

Alternatively, we can express this using standard combination notation:
$$R_x(x) = {}_{2x-2}C_{x-1}$$

This is good, but it can also be written in its factorial form:
$$R_x(x) = \frac{(2x-2)!}{(x-1)!(x-1)!}$$

Finally, if we look at Pascal's Triangle, these are the middle numbers of the triangle:
\begin{center}
\begin{tabular}{c}
\textbf{1} \\
1 \quad 1 \\
1 \quad \textbf{2} \quad 1 \\
1 \quad 3 \quad 3 \quad 1 \\
1 \quad 4 \quad \textbf{6} \quad 4 \quad 1 \\
1 \quad 5 \quad 10 \quad 10 \quad 5 \quad 1 \\
1 \quad 6 \quad 15 \quad \textbf{20} \quad 15 \quad 6 \quad 1 \\
1 \quad 7 \quad 21 \quad 35 \quad 35 \quad 21 \quad 7 \quad 1 \\
1 \quad 8 \quad 28 \quad 56 \quad \textbf{70} \quad 56 \quad 28 \quad 8 \quad 1 \\
1 \quad 9 \quad 36 \quad 84 \quad 126 \quad 126 \quad 84 \quad 36 \quad 9 \quad 1 \\
1 \quad 10 \quad 45 \quad 120 \quad 210 \quad \textbf{252} \quad 210 \quad 120 \quad 45 \quad 10 \quad 1
\end{tabular}
\end{center}

\section{Final Conjecture}

Let us set $R_x(x) = f(x)$, and define the iterative compositions of this function as:
$$f_2(x) = f(f(x))$$
$$f_3(x) = f(f(f(x)))$$

Extending this diagonal iteration to the $x$-th degree, we define:
$$f_x(x) = S(x)$$

This is the \textbf{Seth function}. 

The final conjecture is: Is there a closed-form variation that represents this function? If not, prove that there will never be a closed-form variation for this function, even including future mathematics.

\vspace{1cm}

\begin{center}
\textit{Verified by Mr Ben Wiley, Subject Head of Mathematics at Wynberg Boys' High School.}

\vspace{0.3cm}

\textit{This paper and the final conjecture were developed by Seth Berry (14 years old).}
\end{center}

\end{document}